\documentclass{article}
\usepackage{graphicx} % Required for inserting images

\title{Software-Engineering-Lab}
\author{msaini_be24 }
\date{August 2026}

\documentclass[11pt]{article}
\usepackage[margin=1in]{geometry}
\usepackage{times}
\usepackage{enumitem}
\usepackage{hyperref}
\usepackage{titlesec}
\usepackage{parskip}

\titleformat{\section}{\normalfont\Large\bfseries}{\thesection}{1em}{}
\titleformat{\subsection}{\normalfont\large\bfseries}{\thesubsection}{1em}{}

\title{\textbf{VendorFlow: A Unified SaaS Platform for Local Grocery \& Pharmacy Digitization} \\
\large Project Proposal for UCS503P}
\author{}
\date{}

\begin{document}

\begin{center}
    {\LARGE \textbf{VendorFlow: A Unified SaaS Platform for Local Grocery \& Pharmacy Digitization}}\\[0.3cm]
    {\large Project Proposal for UCS503P}\\[0.3cm]
    Submitted to: Dr. Jeelani\\
    Thapar Institute of Engineering and Technology\\[0.3cm]
    \today
\end{center}

\vspace{0.5cm}

\section{Higher-order goal \dots secondary framing}
This project aims to strengthen local retail resilience by giving small grocery stores and pharmacies
an affordable path to digital commerce. While the broader vision is economic empowerment of
neighborhood businesses, the immediate engineering objective is to translate reliable order
fulfillment and vendor coordination into a measurable, deployable software solution.

\section{Time-to-value \dots primary heuristic}
\begin{itemize}
    \item We will deliver meaningful increments quickly by prototyping the core ordering and
    vendor-routing workflow and validating it with a small set of pilot vendors within a short
    iteration window.
    \item We will prioritize fast feedback over perfection by using weekly iterations (and
    targeting CI-backed automation early) to reduce release friction.
    \item Success is defined by what can be delivered and measured in the first iteration, then
    improved based on vendor and customer feedback.
\end{itemize}

\section{Problem Statement}
Small grocery stores and pharmacies in most Indian cities and towns rely on phone calls,
WhatsApp, spreadsheets, or paper records to manage orders and inventory. Common pain points
include:
\begin{itemize}
    \item \textbf{No online presence:} most vendors cannot afford a custom website, app, or
    delivery infrastructure of their own.
    \item \textbf{Manual, error-prone operations:} inventory and order tracking done by hand
    leads to stockouts, missed orders, and delivery mix-ups.
    \item \textbf{Fragmented discovery:} customers have no single place to find nearby stores,
    compare stock, or track deliveries.
    \item \textbf{Regulatory friction for pharmacies:} prescription-backed orders need a
    verification step that ad-hoc channels (calls, chat apps) do not support reliably.
\end{itemize}

\textbf{Why this matters now (contextual relevance):} Consumer demand for online ordering
continues to grow, but the smallest vendors are the ones least able to build this
infrastructure themselves; a shared platform lets many small businesses digitize at once
without each bearing the full engineering cost.

\section{Proposed Solution \dots Software Proposition}

\subsection{Overview}
Build a full-stack SaaS platform, \textbf{VendorFlow}, with the following components:
\begin{itemize}
    \item \textbf{Customer portal:} discover nearby grocery/pharmacy vendors, browse products,
    upload prescriptions where required, place orders, and track deliveries.
    \item \textbf{Vendor portal:} onboarding and approval workflow, inventory management, order
    management, and a reliability/analytics dashboard.
    \item \textbf{Delivery partner portal:} accept assigned deliveries and update status in
    real time.
    \item \textbf{Admin portal:} vendor approval, prescription verification oversight, platform
    analytics, and audit logs.
    \item \textbf{Smart vendor routing:} route an order to the best-fit vendor using distance,
    live stock, and a computed reliability score.
    \item \textbf{Credit system:} an in-platform credit/wallet mechanism for customers and
    vendors, supporting promotional credits, refunds, and settlement between platform and
    vendors.
\end{itemize}

\subsection{Core workflow}
\begin{enumerate}
    \item Customer browses nearby vendors and adds items to cart (uploading a prescription for
    pharmacy items where required).
    \item System applies smart vendor routing (distance + stock + reliability score) and, for
    pharmacy orders, queues the prescription for manual verification.
    \item Order is confirmed, a delivery partner is assigned, and status updates flow through
    the order lifecycle (Placed $\rightarrow$ Confirmed $\rightarrow$ Out for Delivery
    $\rightarrow$ Delivered).
    \item Customer and vendor receive real-time notifications at each stage; credits/wallet
    balances are updated on completion, refund, or promotional events.
\end{enumerate}

\subsection{Operational constraints for an educational, production-style setting}
\begin{itemize}
    \item Keep the customer- and vendor-facing UIs usable on low-to-mid range devices via
    responsive web design.
    \item Keep the order-matching and routing logic heuristic and explainable rather than
    research-grade; focus on correctness, latency, and auditability.
    \item Design the backend for high throughput, since order placement, routing, and
    notification delivery are on the platform's critical path.
\end{itemize}

\section{Solution Approach \dots Engineering focus}
This is an engineering course project, so we focus on system architecture, performance, and
correctness of the transactional workflow rather than novel algorithm research.

\subsection{Performance-focused backend \dots non-research}
The core transactional services (order placement, vendor routing, inventory updates) are
engineering-heavy and latency-sensitive, so the approach is architectural rather than
algorithmic:
\begin{itemize}
    \item Implement the high-throughput backend services (order intake, routing engine,
    inventory locking) in \textbf{C++} for predictable low-latency performance under load.
    \item Use \textbf{Python} for auxiliary services: analytics jobs, reliability-score
    recomputation, and admin/reporting tooling.
    \item Use \textbf{JavaScript} for the customer, vendor, delivery, and admin front-ends, and
    for lightweight orchestration/API-gateway glue code where appropriate.
    \item Vendor routing itself remains a simple weighted-score heuristic (distance + stock
    availability + reliability score) without claiming state-of-the-art optimization.
\end{itemize}

\subsection{System architecture \dots web-based solution}
A multi-tier architecture:
\begin{itemize}
    \item \textbf{Frontend:} responsive web apps for the four portals (customer, vendor,
    delivery partner, admin), built in JavaScript.
    \item \textbf{High-performance backend (C++):} REST/WebSocket services for authentication,
    order lifecycle, vendor routing, inventory locking, and the credit/wallet ledger.
    \item \textbf{Auxiliary services (Python):} background jobs (reliability score updates,
    routing-timeout handling, notification queue processing) and analytics pipelines.
    \item \textbf{Cache/queue layer:} \textbf{Redis} for caching, rate limiting, session data,
    and background job/notification queues.
    \item \textbf{Primary database:} \textbf{PostgreSQL} for users, vendors, products,
    inventory, orders, credit/wallet ledger, reviews, and audit logs.
    \item \textbf{Real-time layer:} Socket.IO (or an equivalent WebSocket layer) for real-time
    order and delivery status notifications.
\end{itemize}

\subsection{CI/CD and fast delivery enabler}
CI/CD will be used to:
\begin{itemize}
    \item Automatically build and run tests (C++ unit tests, Python service tests, frontend
    tests) on every change.
    \item Produce repeatable deployment artifacts to staging for each service.
    \item Enable safe and frequent releases of incremental features across portals.
\end{itemize}

\section{Evaluation Criterion \dots measurable and attributable}
We define success using measurable metrics tied to the core order-fulfillment workflow.

\subsection{Primary evaluation metric}
\textbf{Order-to-confirmation latency (OCL):} time between order placement and vendor
confirmation (post-routing) in the system.
\begin{itemize}
    \item Target: median OCL $\leq$ 60 seconds during pilot window.
    \item Attribution: measured from backend timestamps (order-placed event vs.
    vendor-confirmed event).
\end{itemize}

\subsection{Secondary evaluation metrics}
\begin{itemize}
    \item \textbf{Delivery success rate:} percentage of orders marked Delivered without
    escalation or cancellation.
    \item \textbf{Vendor reliability accuracy:} correlation between computed reliability score
    and actual on-time fulfillment during the pilot.
    \item \textbf{Prescription verification turnaround:} median time from upload to
    admin-verified status for pharmacy orders.
    \item \textbf{System reliability:} availability of staging/production for login, ordering,
    and routing during business hours (e.g., $\geq$ 99\% uptime in pilot).
    \item \textbf{Backend throughput:} sustained requests/second the C++ order and routing
    services can serve under load-test conditions.
\end{itemize}

\subsection{Pilot validation plan \dots high-level}
\begin{itemize}
    \item Recruit 10--20 pilot customers and 3--5 vendor accounts (grocery and pharmacy mix,
    simulated where real vendors are unavailable).
    \item Compare metrics before/after within the iteration window, using short interviews or
    existing process observations to estimate baseline manual-process timings.
\end{itemize}

\section{Scalability \dots with foresight}
The proposition should scale in both theoretical and practical senses.
\begin{enumerate}
    \item \textbf{Scalable (theoretically):} The system should support growth in vendors,
    orders, and concurrent portal usage by:
    \begin{itemize}
        \item decoupling the C++ backend, Python auxiliary services, and frontends,
        \item enabling pagination, indexing, and read replicas for common PostgreSQL queries,
        \item using Redis for caching and rate limiting to shield the primary database,
        \item designing backend APIs to be stateless where possible.
    \end{itemize}
    \item \textbf{Operationally deployable (practically):} The system should run on common
    infrastructure:
    \begin{itemize}
        \item managed PostgreSQL and Redis instances,
        \item containerized deployment for the C++ backend, Python services, and frontends,
        \item CI/CD pipeline for staging/production across all services.
    \end{itemize}
\end{enumerate}

\section{Engine availability heuristic}
A mature set of ``engines'' (tooling/services) is available to implement this as a
production-style project:
\begin{itemize}
    \item High-performance C++ web/service framework for the core backend.
    \item Python framework for auxiliary/background services and analytics.
    \item JavaScript framework for the four frontend portals.
    \item PostgreSQL as the managed relational database.
    \item Redis for caching, rate limiting, and job/notification queues.
    \item Authentication approach (JWT-based, with role-based access control).
    \item Object storage for prescription images and product photos.
    \item Test frameworks for unit/integration tests across C++, Python, and JavaScript
    components.
    \item CI runner and staging deployment target.
\end{itemize}
We will use these mature components to focus on workflow design, backend performance, and
measurement rather than inventing new infrastructure.

\section{Project scope and deliverables \dots iteration-friendly}

\subsection{Initial deliverable \dots first iteration}
\begin{itemize}
    \item Customer, Vendor, Delivery Partner, and Admin portals with JWT authentication and
    role-based access control.
    \item Vendor onboarding \& approval workflow; product \& inventory management.
    \item Smart vendor routing (distance + stock + reliability) implemented in the C++ backend.
    \item Shopping cart, checkout, and order lifecycle management.
    \item Prescription upload with manual admin verification.
    \item Delivery assignment and status-based tracking.
    \item Basic logging and timestamps for OCL measurement.
    \item CI: build + backend (C++/Python) unit tests + linting.
    \item CD to staging with a single pipeline job.
\end{itemize}

\subsection{Subsequent deliverables}
\begin{itemize}
    \item Credit/wallet system for customers and vendors (promotional credits, refunds,
    settlement).
    \item Real-time notifications via Socket.IO, plus email and in-app notifications.
    \item Reviews \& ratings, feeding into the vendor reliability score.
    \item Vendor \& admin analytics dashboards.
    \item Redis-backed caching and rate limiting hardened for scale.
    \item Background jobs for notification queueing, routing timeouts, and reliability-score
    recomputation.
    \item Audit logs and activity history across vendor and admin actions.
    \item Improved search/filtering and indexed queries for scale readiness.
\end{itemize}

\section{Risks and mitigations}
\begin{itemize}
    \item \textbf{Data privacy / consent:} minimize collected data, especially prescription
    images; store only what is needed and restrict access to verified admins.
    \item \textbf{Vendor adoption:} keep the vendor portal simple, with minimal steps for
    onboarding, inventory updates, and order confirmation.
    \item \textbf{Backend complexity (C++):} mitigate with strong test coverage, clear service
    boundaries, and CI-enforced builds before merge.
    \item \textbf{Evaluation measurement ambiguity:} instrument timestamps and define clear
    order-lifecycle states before development begins.
    \item \textbf{Operational issues in pilot:} use staging early; ensure CI/CD produces
    repeatable deployments for all services.
\end{itemize}

\section{Summary}
This project proposes VendorFlow, a deployable SaaS platform that helps small grocery stores
and pharmacies digitize inventory, ordering, delivery, and vendor management, with a
high-performance C++ backend, Python auxiliary services, PostgreSQL and Redis for storage and
caching, and JavaScript frontends across four portals. It meets the course criteria by
emphasizing time-to-value through rapid iterations, implementing CI/CD to support frequent safe
releases, and using measurable evaluation criteria (especially order-to-confirmation latency).
It remains focused on engineering workflow, backend performance, and scalability readiness
rather than research-driven novelty.

\end{document}